% Fragmento público generado desde propietarios íntegros.
% Residencia: 52.8.5.
% Fuente: ../incoming/radion_monografia_20260827/manuscrito/sections/09_sector_90120.tex:55-114
\subsubsection{Publication spectrale de la queue d'ordres \(j\ge2\)}

\paragraph{Définition de la fonctionnelle spectrale}

Pour \(m\in\{90,120\}\), soit
\begin{equation}
L_m^{(2)}(q)=\sum_{j\ge2}\frac{q^{mj}}{j}
=-\log(1-q^m)-q^m.
\label{eq:L2}
\end{equation}
La lecture d'une paire est
\begin{equation}
T_h^{(2)}=\frac{180}{\pi}
\Bigl[
L_{90}^{(2)}(q_{h,-})-L_{90}^{(2)}(q_{h,+})
-L_{120}^{(2)}(q_{h,-})+L_{120}^{(2)}(q_{h,+})
\Bigr].
\label{eq:T-h}
\end{equation}
La reconstruction des six paires donne
\begin{equation}
F_{\mathrm{spec}}=6T_h^{(2)}
=5.1499235153094574410\ldots\times10^{-8}\,{}^\circ.
\label{eq:Fspec}
\end{equation}

La retenue du radion est
\[
F_{\mathrm{acarreo}}=\epsR^\circ
=5.1383016758575282072\ldots\times10^{-8}\,{}^\circ.
\]
Ils ne coïncident pas :
\begin{equation}
\chi_R:=F_{\mathrm{spec}}-F_{\mathrm{acarreo}}
=1.16218394519292338\ldots\times10^{-10}\,{}^\circ.
\label{eq:chiR}
\end{equation}

L’égalité
\[
F_{\mathrm{acarreo}}=F_{\mathrm{spec}}-\chi_R
\]
est un changement de carte valable une fois les deux sorties construites. Il ne doit pas
être présenté comme génération indépendante de \(\epsR^\circ\) si \(\chi_R\) est
défini par la soustraction elle-même. La queue est clôturée et significative ; ce qui est faux,
c'est de l'identifier directement à la retenue.

\begin{falsifier}
Après conversion de la queue exprimée en degrés dans l'unité angulaire interne,
\[
\frac{\pi}{180^\circ}F_{\mathrm{spec}}-\epsone
=2.0283936357433839\ldots\times10^{-12}\ne0.
\]
Par conséquent, ni le log complet ni la queue \(j\ge2\), même après application de la
carte d'unités, ne sont \(\epsone\). Une
seconde factorisation spectrale autonome nécessiterait de générer \(\chi_R\) sans
utiliser la différence comme définition ; la clôture APP--TPK du radion ne dépend pas
de cette promotion.
\end{falsifier}

